\documentclass[
    language=english,
    doctype=article,
    institution=none,
]{../../omnilatex}

\RequirePackage{config/document-settings}
\addbibresource{../../bib/bibliography.bib}

\begin{document}

\maketitle

\begin{abstract}
This document specifies the CloudVault REST API for programmatic access to
cloud object storage.  All endpoints return JSON responses and accept JSON
request bodies unless otherwise noted.  The base URL for all requests is
\texttt{https://api.cloudvault.io/v2}.
\end{abstract}

\tableofcontents

\section{Authentication}

All API requests require authentication via an API key or OAuth 2.0 bearer
token.  Include the token in the \texttt{Authorization} header:

\begin{verbatim}
Authorization: Bearer <your-api-token>
\end{verbatim}

\subsection{Obtaining API Keys}

Navigate to \texttt{Settings > API Keys} in the CloudVault dashboard to
generate a key pair.  The secret key is shown only once at creation time.
Keys are scoped to a single project and can be revoked at any time.

\subsection{OAuth 2.0 Flow}

For user-delegated access, implement the OAuth 2.0 authorization code flow:

\begin{enumerate}
    \item Redirect the user to \texttt{https://auth.cloudvault.io/authorize}
          with \texttt{client\_id}, \texttt{redirect\_uri}, and
          \texttt{scope} parameters.
    \item Exchange the authorization code for an access token via
          \texttt{POST /oauth/token}.
    \item Use the access token in subsequent API requests.
\end{enumerate}

\section{Endpoints}

\subsection{Buckets}

\begin{table}[htbp]
    \centering
    \caption{Bucket management endpoints.}
    \label{tab:buckets}
    \begin{tabular}{@{}llp{5cm}@{}}
        \toprule
        \textbf{Method} & \textbf{Endpoint} & \textbf{Description} \\
        \midrule
        \texttt{GET}    & \texttt{/buckets}              & List all buckets
            in the project \\
        \texttt{POST}   & \texttt{/buckets}              & Create a new bucket \\
        \texttt{GET}    & \texttt{/buckets/\{name\}}     & Get bucket metadata
            and configuration \\
        \texttt{DELETE} & \texttt{/buckets/\{name\}}     & Delete a bucket
            (must be empty) \\
        \texttt{PUT}    & \texttt{/buckets/\{name\}/policy}
            & Update bucket access policy \\
        \bottomrule
    \end{tabular}
\end{table}

\subsection{Objects}

\begin{table}[htbp]
    \centering
    \caption{Object storage endpoints.}
    \label{tab:objects}
    \begin{tabular}{@{}llp{5cm}@{}}
        \toprule
        \textbf{Method} & \textbf{Endpoint} & \textbf{Description} \\
        \midrule
        \texttt{GET}    & \texttt{/buckets/\{name\}/objects}
            & List objects with optional prefix filter \\
        \texttt{PUT}    & \texttt{/buckets/\{name\}/objects/\{key\}}
            & Upload or overwrite an object \\
        \texttt{GET}    & \texttt{/buckets/\{name\}/objects/\{key\}}
            & Download an object \\
        \texttt{HEAD}   & \texttt{/buckets/\{name\}/objects/\{key\}}
            & Retrieve object metadata without body \\
        \texttt{DELETE} & \texttt{/buckets/\{name\}/objects/\{key\}}
            & Permanently delete an object \\
        \texttt{POST}   & \texttt{/buckets/\{name\}/objects/\{key\}/copy}
            & Copy an object within or across buckets \\
        \bottomrule
    \end{tabular}
\end{table}

\section{Request and Response Examples}

\subsection{Create a Bucket}

\textbf{Request:}

\begin{verbatim}
POST /v2/buckets
Content-Type: application/json

{
    "name": "analytics-logs-2026",
    "region": "us-east-1",
    "storage_class": "standard",
    "versioning_enabled": true
}
\end{verbatim}

\textbf{Response (201 Created):}

\begin{verbatim}
{
    "name": "analytics-logs-2026",
    "region": "us-east-1",
    "storage_class": "standard",
    "versioning_enabled": true,
    "created_at": "2026-06-12T14:30:00Z",
    "arn": "arn:cloudvault:us-east-1:48291:bucket/analytics-logs-2026"
}
\end{verbatim}

\subsection{Upload an Object}

\textbf{Request:}

\begin{verbatim}
PUT /v2/buckets/analytics-logs-2026/objects/events/2026-06-12.jsonl
Content-Type: application/octet-stream
Content-MD5: 1B2M2Y8AsgTpgAmY7PhCfg==

<binary or streaming body>
\end{verbatim}

\textbf{Response (200 OK):}

\begin{verbatim}
{
    "key": "events/2026-06-12.jsonl",
    "size": 104857600,
    "etag": "d41d8cd98f00b204e9800998ecf8427e",
    "content_type": "application/octet-stream",
    "last_modified": "2026-06-12T14:35:22Z",
    "storage_class": "standard"
}
\end{verbatim}

\subsection{List Objects with Pagination}

\textbf{Request:}

\begin{verbatim}
GET /v2/buckets/analytics-logs-2026/objects?prefix=events/&max_keys=100
    &continuation_token=abc123
\end{verbatim}

\textbf{Response (200 OK):}

\begin{verbatim}
{
    "objects": [
        {"key": "events/2026-06-12.jsonl", "size": 104857600,
         "last_modified": "2026-06-12T14:35:22Z"},
        {"key": "events/2026-06-11.jsonl", "size": 98765432,
         "last_modified": "2026-06-11T23:59:58Z"}
    ],
    "is_truncated": true,
    "next_continuation_token": "def456",
    "total_count": 1247
}
\end{verbatim}

\section{Error Codes}

The API uses standard HTTP status codes.  Error responses include a JSON body
with a machine-readable error code and a human-readable message.

\begin{table}[htbp]
    \centering
    \caption{API error codes.}
    \label{tab:errors}
    \begin{tabular}{@{}llp{5cm}@{}}
        \toprule
        \textbf{HTTP} & \textbf{Error Code} & \textbf{Description} \\
        \midrule
        400 & \texttt{INVALID\_REQUEST}    & The request body or parameters are
            malformed or missing required fields. \\
        401 & \texttt{UNAUTHORIZED}         & The API key or token is missing,
            expired, or invalid. \\
        403 & \texttt{FORBIDDEN}            & The authenticated principal lacks
            permission for the requested operation. \\
        404 & \texttt{NOT\_FOUND}           & The requested bucket or object
            does not exist. \\
        409 & \texttt{CONFLICT}             & The bucket already exists or a
            concurrent modification was detected. \\
        429 & \texttt{RATE\_LIMITED}        & Too many requests.  Retry after
            the duration specified in the \texttt{Retry-After} header. \\
        500 & \texttt{INTERNAL\_ERROR}      & An unexpected server error
            occurred.  Contact support with the request ID. \\
        503 & \texttt{SERVICE\_UNAVAILABLE} & The service is temporarily
            unavailable.  Retry with exponential backoff. \\
        \bottomrule
    \end{tabular}
\end{table}

Example error response:

\begin{verbatim}
{
    "error": {
        "code": "NOT_FOUND",
        "message": "Bucket 'my-bucket' does not exist.",
        "request_id": "req-7f3a9b2c"
    }
}
\end{verbatim}

\section{Rate Limiting}

CloudVault enforces rate limits per API key to ensure fair usage:

\begin{table}[htbp]
    \centering
    \caption{Rate limit tiers.}
    \label{tab:ratelimits}
    \begin{tabular}{@{}lrr@{}}
        \toprule
        \textbf{Tier} & \textbf{Requests / Minute} & \textbf{Burst} \\
        \midrule
        Free        & 60   & 10  \\
        Pro         & 600  & 100 \\
        Enterprise  & 6\,000 & 1\,000 \\
        \bottomrule
    \end{tabular}
\end{table}

Rate limit headers are included in every response:

\begin{verbatim}
X-RateLimit-Limit: 600
X-RateLimit-Remaining: 587
X-RateLimit-Reset: 1718212800
\end{verbatim}

When the limit is exceeded, the API returns \texttt{429 Too Many Requests}.
Implement exponential backoff with jitter to handle rate limiting gracefully.

\section{Versioning}

The API is versioned via the URL path (\texttt{/v2}).  When a new major version
is released, the previous version enters a 12-month deprecation period.  Deprecation
headers are included in responses for deprecated endpoints.

\begin{verbatim}
Sunset: Sat, 12 Jun 2027 00:00:00 GMT
Deprecation: true
Link: <https://docs.cloudvault.io/api/v3>; rel="successor-version"
\end{verbatim}

\printbibliography[heading=none]

\end{document}
